\section{Published model and notation}
\label{sec:model}

We work with the model exactly as specified by \cite{ShyStenbacka2005}, except where an auxiliary restriction is introduced explicitly and labeled as such.  A firm uses a continuum of inputs of total measure \(\phi\).  Under the constant outsourcing-cost advantage used in the paper's equilibrium analysis, firm \(j\) chooses the measure \(i_j\) of outsourced inputs subject to
\begin{equation}
0\le i_j\le \phi .
\label{eq:sourcing-domain}
\end{equation}
Its final-good marginal cost is
\begin{equation}
c_j
=
H(\phi-i_j)+\frac{\gamma\phi^2}{2}
=
C_0-Hi_j,
\qquad
C_0\equiv H\phi+\frac{\gamma\phi^2}{2},
\label{eq:mc}
\end{equation}
and the monitoring cost is
\begin{equation}
M(i_j)=i_j^2.
\label{eq:monitoring}
\end{equation}
Thus a larger \(i_j\) lowers marginal production cost but raises a fixed monitoring cost.  Define
\begin{equation}
D\equiv a-C_0>0.
\label{eq:D}
\end{equation}
No free coefficient multiplying \(i_j^2\) is introduced in the baseline model.

\subsection{Cournot product-market competition}

There are \(N\ge2\) final-good firms.  Stage I consists of the simultaneous sourcing choices \((i_1,\ldots,i_N)\in[0,\phi]^N\).  In Stage II the firms choose nonnegative quantities,
\[
q_j\ge0,
\]
under inverse demand
\begin{equation}
p=a-bQ,\qquad Q=\sum_{k=1}^N q_k .
\label{eq:inverse-demand}
\end{equation}
Firm \(j\)'s payoff is
\begin{equation}
\Pi_j=(p-c_j)q_j-i_j^2 .
\label{eq:cournot-profit}
\end{equation}
The maintained second-order condition used in the source is
\begin{equation}
b>\left(\frac{HN}{N+1}\right)^2.
\label{eq:source-soc}
\end{equation}
Our equilibrium analysis keeps the primitive constraints \eqref{eq:sourcing-domain} and \(q_j\ge0\) at every history.  We focus on pure Stage-I strategies and the unique pure quantity continuation established below.  For general \(N\), we characterize the unique symmetric pure Stage-I action; we do not claim uniqueness of the complete Stage-I equilibrium set.

\subsection{Hotelling product-market competition}

The paper also studies a differentiated duopoly with firms located at the endpoints of a unit interval, market density \(n\), and transport parameter \(\tau>0\).  After sourcing choices \((i_A,i_B)\), consumers compare
\[
u_A(z)=\omega-p_A-\tau z,
\qquad
u_B(z)=\omega-p_B-\tau(1-z).
\]
Under full market coverage, firm A's actual market share is
\begin{equation}
s_A(p_A,p_B)
=
\min\left\{
1,\,
\max\left\{
0,\,
\frac12+\frac{p_B-p_A}{2\tau}
\right\}
\right\}.
\label{eq:clipped-share}
\end{equation}
The literal source game places no lower bound \(p_j\ge c_j\) on prices.  We therefore analyze its pure-price equilibrium correspondence directly.  Later, as a separate robustness exercise, we study an auxiliary no-loss price game that explicitly restricts \(p_j\ge c_j\).  That auxiliary restriction is not attributed to \cite{ShyStenbacka2005} and is not treated as an equilibrium refinement generated by elimination of weakly dominated strategies.
